% Options for packages loaded elsewhere
\PassOptionsToPackage{unicode}{hyperref}
\PassOptionsToPackage{hyphens}{url}
%
\documentclass[
  10pt,
  runningheads]{llncs}
\usepackage{amsmath,amssymb}
\usepackage{iftex}
\ifPDFTeX
  \usepackage[T1]{fontenc}
  \usepackage[utf8]{inputenc}
  \usepackage{textcomp} % provide euro and other symbols
\else % if luatex or xetex
  \usepackage{unicode-math} % this also loads fontspec
  \defaultfontfeatures{Scale=MatchLowercase}
  \defaultfontfeatures[\rmfamily]{Ligatures=TeX,Scale=1}
\fi
\usepackage{lmodern}
\ifPDFTeX\else
  % xetex/luatex font selection
\fi
% Use upquote if available, for straight quotes in verbatim environments
\IfFileExists{upquote.sty}{\usepackage{upquote}}{}
\IfFileExists{microtype.sty}{% use microtype if available
  \usepackage[]{microtype}
  \UseMicrotypeSet[protrusion]{basicmath} % disable protrusion for tt fonts
}{}
\makeatletter
\@ifundefined{KOMAClassName}{% if non-KOMA class
  \IfFileExists{parskip.sty}{%
    \usepackage{parskip}
  }{% else
    \setlength{\parindent}{0pt}
    \setlength{\parskip}{6pt plus 2pt minus 1pt}}
}{% if KOMA class
  \KOMAoptions{parskip=half}}
\makeatother
\usepackage{xcolor}
\usepackage{longtable,booktabs,array}
\usepackage{calc} % for calculating minipage widths
% Correct order of tables after \paragraph or \subparagraph
\usepackage{etoolbox}
\makeatletter
\patchcmd\longtable{\par}{\if@noskipsec\mbox{}\fi\par}{}{}
\makeatother
% Allow footnotes in longtable head/foot
\IfFileExists{footnotehyper.sty}{\usepackage{footnotehyper}}{\usepackage{footnote}}
\makesavenoteenv{longtable}
\setlength{\emergencystretch}{3em} % prevent overfull lines
\providecommand{\tightlist}{%
  \setlength{\itemsep}{0pt}\setlength{\parskip}{0pt}}
\setcounter{secnumdepth}{-\maxdimen} % remove section numbering
\ifLuaTeX
  \usepackage{selnolig}  % disable illegal ligatures
\fi
\usepackage{bookmark}
\IfFileExists{xurl.sty}{\usepackage{xurl}}{} % add URL line breaks if available
\urlstyle{same}
\hypersetup{
  hidelinks,
  pdfcreator={LaTeX via pandoc}}

\title{Evidence--Rating Alignment for Evidence-Grounded HRMS Performance Evaluation}
\titlerunning{Evidence--Rating Alignment in HRMS}
\author{Divesh Singh\\
GLA University, Mathura, India\\
\email{divesh.singh\_phd.cs24@gla.ac.in}\\[8pt]
Dr. Dilip Kumar Sharma\\
GLA University, Mathura, India\\
\email{dilip.sharma@gla.ac.in}}
\authorrunning{D. Singh and D. K. Sharma}
\institute{}
\date{}

\begin{document}
\setlength{\parskip}{0pt}
\setlength{\parindent}{1em}
\maketitle
\begin{abstract}

Employee-performance evaluation in human resource management systems
(HRMS) can produce ratings and explanations that appear consistent while
lacking criterion-specific evidential support. Agreement between an
evaluator and a manager is not, by itself, evidence that the rating is
justified. This paper specifies Evidence--Rating Alignment (ERA), a
proposed measurement layer for auditing the relationship between role
criteria, available evidence, ratings, and explanation claims. ERA
combines Evidence Coverage (EC), Evidence Relevance (ER), Criterion
Agreement (CA), Evidence Traceability (ET), and a Contradiction Penalty
(CP): ERA = 0.25EC + 0.25ER + 0.25CA + 0.15ET + 0.10(1-CP). The paper
evaluates the arithmetic and data structure of an existing synthetic
pilot workbook containing 100 constructed cases across 10 job roles and
500 criterion-level annotation rows. All 100 reference records remain
explicitly provisional, and all core annotation fields are blank;
consequently, they cannot serve as independently validated ground truth.
Recalculation of the stored ERA composite produced a maximum absolute
discrepancy of 0.0006 and no row exceeded a 0.001 tolerance. An
exploratory weight-sensitivity analysis on the synthetic component
values showed greater rank stability under local perturbations around
the proposed weights than under broad perturbations across the weight
simplex. These are computational and mathematical audit results only.
They do not demonstrate AI accuracy, human agreement,
employee-performance validity, fairness, explanation usefulness, or
real-world effectiveness. ERA is presented as a transparent, testable
design proposal whose validity remains an open empirical question.

\keywords{Explainable AI \and Human resource management systems \and Employee performance evaluation \and Evidence grounding \and Auditability \and Synthetic data \and Composite indicators}
\end{abstract}

\section{Introduction}\label{introduction}

AI-supported HRM can help organize records and support decisions, but
employee performance is context-dependent and consequential. Research on
AI in HRM identifies concerns involving data quality, accountability,
fairness, privacy, employee experience, and the role of human judgment
{[}1,2,11,12{]}. Research on explainable AI (XAI) also cautions that an
explanation should be evaluated in relation to its purpose and audience
rather than assumed trustworthy merely because it is readable {[}3,4{]}.

Consider a hypothetical employee for whom a manager and an AI system
both assign 4/5. The agreement may look reassuring, but both assessments
may have omitted a missed deadline or relied on a metric unrelated to
the employee's role. Agreement alone cannot reveal whether the rating is
supported by appropriate records. This example is illustrative, not an
observed organizational case.

This paper addresses the measurement gap between a performance
criterion, the evidence available and used, the rating assigned, and the
claims offered to justify it. It introduces Evidence--Rating Alignment
(ERA) as a proposed audit layer rather than a predictive model or a
complete definition of fairness. ERA is intended to make evidence
coverage, relevance, reference agreement, traceability, and
contradiction visible in a common record. The central contribution is a
transparent operational specification and an audit of the currently
available synthetic pilot dataset. No new human annotations, real
employee records, or fresh LLM outputs are introduced.

\section{Research Problem, Objectives, and
Contributions}\label{research-problem-objectives-and-contributions}

Research problem: how can an HRMS represent whether a criterion-level
performance rating is supported by role-relevant evidence, make
contradictions and missing support visible, and preserve a traceable
record for review and contestation?

The paper has four objectives: (1) define the five ERA components and a
normalized composite; (2) specify how the composite can be embedded in
an evidence-to-rating workflow; (3) audit the structure and arithmetic
of the existing 100-case synthetic workbook; and (4) examine how
rankings from its provisional component values respond to alternative
weights. These objectives are intentionally narrower than evaluating a
live AI system or determining whether ERA improves human decisions.

The contributions are: (a) a criterion-level operational definition of
evidence--rating alignment; (b) a proposed audit record linking
criteria, evidence identifiers, ratings, contradictions, and review
actions; (c) a reproducible dataset and formula audit; and (d) an
exploratory weight-sensitivity analysis. The paper does not claim that
the composite is validated, that the synthetic reference scores are
independent ground truth, or that ERA outperforms existing approaches.

\section{Related Work and Research
Gap}\label{related-work-and-research-gap}

Responsible-AI research in HRM has examined accountability, privacy,
fairness, transparency, and employee involvement
{[}1,2,5,11,12,20,23{]}. These works motivate governance-aware designs,
but broad principles do not themselves specify a case-by-criterion
measure of whether a particular rating is grounded in particular
evidence. Reviews of algorithmic decision-making in employment also show
that fairness and discrimination concerns arise across recruitment and
development settings {[}5,20{]}.

XAI research offers methods for explaining predictions, including local
explanation techniques such as LIME {[}3,4{]}. However, explanation
fluency or feature attribution is not equivalent to verifying that a
record supports a criterion-specific claim. ERA therefore treats
evidence relevance and traceability as separate dimensions, rather than
treating the presence of an explanation as proof of evidential support.

Research on algorithmic management and organizational justice emphasizes
worker experience, power, procedural fairness, dignity, and
contestability {[}6--9,13,14,17,18{]}. A numerical evidence-alignment
score cannot replace these dimensions. A high ERA score should not be
interpreted as proof of fairness, legitimacy, or acceptable employment
procedure.

A directly related study by Li, Zhou, and Xu examines LLM-based
performance evaluation against an external expert-consensus reference
{[}10{]}. That design illustrates the importance of independent
benchmarks when evaluating ratings. The present paper does not run an
LLM and has no independently adjudicated benchmark. Its contribution is
complementary and methodological: it specifies a candidate
evidence-alignment layer and audits the arithmetic of a synthetic
prototype. The claimed gap is limited to the selected sources discussed
here; it is not a claim that no related method exists anywhere in the
literature.

\section{Literature Search and Evidence
Boundary}\label{literature-search-and-evidence-boundary}

The source base is a targeted, structured search for a focused narrative synthesis, not a systematic review. The source manuscript records 11 October 2026 as the search date. Topics included responsible AI in HRM, employee-performance evaluation, explainable AI, algorithmic fairness, human oversight, and benchmark validity. Eligibility focused on English-language scholarly work and authoritative guidance directly relevant to these topics. The search was single-researcher and targeted; independent duplicate screening, a complete deduplication log, formal risk-of-bias assessment, and PRISMA flow counts were not preserved. Therefore, the paper does not claim systematic-review status, comprehensive coverage, or a verified 43-study corpus. Reference [23] is treated only as an external published source.

\section{ERA Framework and Mathematical
Specification}\label{era-framework-and-mathematical-specification}

For case i, let k index predefined role criteria and j index evidence
items. Criterion weights wᵢₖ are non-negative and sum to one over the
criteria assessable for that case. Let rᵢ be the evaluated rating and bᵢ
the reference rating on the same anchored 1--5 scale. The five ERA
components are normalized to {[}0,1{]}. A component is undefined if its
denominator is zero; in that situation, the case should be flagged and a
complete composite should not be silently manufactured.

The proposed composite is:

\textbf{ERAᵢ = 0.25 ECᵢ + 0.25 ERᵢ + 0.25 CAᵢ + 0.15 ETᵢ + 0.10 (1 -
CPᵢ)}

\subsection{Evidence Coverage (EC)}\label{evidence-coverage-ec}

Let Aᵢₖ = 1 when criterion k is assessable and 0 otherwise, and Rᵢₖ = 1
when at least one required, relevant evidence item was retrieved for
that criterion. Then ECᵢ = Σₖ wᵢₖ Aᵢₖ Rᵢₖ / Σₖ wᵢₖ Aᵢₖ. The denominator
must be positive. Coverage measures whether the required evidence was
available to the evaluator; it does not establish that the evidence is
accurate.

\subsection{Evidence Relevance (ER)}\label{evidence-relevance-er}

For each evidence item actually used, a pre-registered rubric assigns
sᵢⱼ ∈ \{1,2,3,4,5\}; normalize qᵢⱼ = (sᵢⱼ - 1)/4. With non-negative use
weights vᵢⱼ, ERᵢ = Σⱼ vᵢⱼ qᵢⱼ / Σⱼ vᵢⱼ. Equal weights are the default
unless another rule is justified before analysis. Evidence that was
available but not used is excluded from this component and should be
recorded separately as a retrieval omission. If no item was used, ER is
undefined.

\subsection{Criterion Agreement (CA)}\label{criterion-agreement-ca}

CAᵢ = max(0, 1 - \textbar rᵢ - bᵢ\textbar/4), where rᵢ and bᵢ use the
same anchored 1--5 scale. CA expresses agreement with the selected
reference, not proof that the reference is correct. Because the current
reference values are synthetic and provisional, CA values calculated
from them must not be interpreted as agreement with validated
employee-performance ground truth.

\subsection{Contradiction Penalty
(CP)}\label{contradiction-penalty-cp}

Let cᵢⱼ = 1 when evidence item j materially contradicts the evaluated
claim under a defined coding rubric and 0 otherwise. CPᵢ = Σⱼ vᵢⱼ cᵢⱼ /
Σⱼ vᵢⱼ over assessable evidence items used. Missing evidence is not
automatically contradictory. If no assessable evidence was used, CP is
undefined.

\subsection{Evidence Traceability
(ET)}\label{evidence-traceability-et}

Let Lᵢ be the number of atomic rating or explanation claims that require
support and Tᵢ the number that include a valid evidence identifier and a
verifiable claim-to-record mapping. ETᵢ = Tᵢ/Lᵢ. A citation or
identifier alone is insufficient: the record must exist, be accessible
to an authorized reviewer, and support the relevant claim. If Lᵢ = 0, ET
is undefined.

\subsection{Weighting and
Interpretation}\label{weighting-and-interpretation}

The weights sum to one and are a transparent, theory-led starting
configuration, not empirically estimated parameters. EC, ER, and CA each
receive 0.25; ET receives 0.15; and CP contributes a 0.10 penalty term.
These choices express design priorities rather than proven relative
importance. Composite-indicator research recommends explicit weight
assumptions and sensitivity checks {[}21,22{]}. The additive score is
compensatory, so it must be shown alongside component values and flags.
A high score must not override missing evidence, material contradiction,
or the need for human review, and it must not automatically determine
pay, promotion, termination, or other employment outcomes.

\section{Proposed XAI--HRMS Workflow and Audit
Record}\label{proposed-xaihrms-workflow-and-audit-record}

The proposed workflow begins with authorized records and a frozen
role-specific criterion rubric. A retrieval or evidence-selection stage
links each assessable criterion to evidence identifiers, provenance,
source type, and time window. A rating stage records the criterion-level
score and the evidence actually used. An audit stage computes the ERA
components, records missing and contradictory evidence, and identifies
unsupported claims. A human-review stage permits the authorized reviewer
to inspect the records, override or abstain from a rating, record a
rationale, and initiate correction or appeal. The architecture is a
design specification, not a deployed HRMS service.

For each criterion-level output, the audit record should include: case
and criterion IDs; rubric version; rating and scale; evidence IDs used;
evidence provenance and timestamps; evidence relevance and validity
coding; missing or contradictory evidence flags; atomic explanation
claims and their evidence links; component values; composite ERA;
undefined-component flags; reviewer action and rationale; and a
versioned audit timestamp. The record supports traceability and
contestability but does not itself prove that a rating is fair or valid.

\section{Explanation, Fairness, and Human Oversight
Measures}\label{explanation-fairness-and-human-oversight-measures}

The current dataset contains no generated natural-language explanations
that were independently reviewed. Therefore explanation quality is
specified, not measured. For a future evaluation, explanations should be
decomposed into atomic factual claims. Let N be the number of eligible
claims, S the number fully supported, U the number unsupported or
contradicted, V the number unverifiable because required records are
unavailable, and P the number partially supported. These categories
should be mutually exclusive and sum to N. Evidence Support Rate is S/N;
Unsupported Claim Rate is U/N; Unverifiable Claim Rate is V/N; and
Partial Support Rate is P/N. If N=0, these metrics are undefined.

Evidence identification precision and recall can be calculated against a
separately frozen evidence key: EIP =
\textbar R∩G\textbar/\textbar R\textbar{} and EIR =
\textbar R∩G\textbar/\textbar G\textbar, where R is the set of returned
evidence IDs and G is the set of required IDs. Their harmonic mean is
reported only when defined. For planted contradictions, error-detection
precision, recall, F1, and false-positive rate can be measured. These
metrics are not reported as results here because no qualifying model
outputs or independently scored explanations are present.

Explanation usefulness should be measured through separately reported
reviewer ratings of clarity, evidence sufficiency, actionability, and
contestability. Fairness requires a separate analysis of decision
processes and error patterns; ERA itself contains no protected-attribute
or subgroup fairness term. Sensitive attributes should be used only
where lawfully authorized and ethically approved. Meaningful human
oversight requires the ability to inspect evidence, challenge a rating,
correct records, abstain, and appeal. These are design requirements, not
evaluated outcomes of the current study.

\section{Data and Methods}\label{data-and-methods}

\subsection{Available Dataset}\label{available-dataset}

The analysis uses the existing workbook
ERA\_Pilot\_Dataset\_100\_Cases\_QC\_Corrected.xlsx, sheet
``100\_Cases\_QC\_Corrected'', plus its role-criteria and
annotation-form sheets. The case table contains 100 synthetic records,
100 unique case IDs, and 10 roles with 10 cases per role. The designed
case-pattern distribution is: aligned (20), ai\_human\_agree\_wrong
(20), contradiction (15), strong (15), weak (15), mixed (10), and
missing (5). These proportions are features of the constructed dataset
and must not be interpreted as the prevalence of these patterns in real
organizations.

The workbook also contains 500 criterion-level annotation rows (five per
case). The core fields for evidence relevance, evidence validity,
criterion benchmark rating, and contradiction coding are blank in all
500 rows. Every case's benchmark status explicitly states that human
annotation/adjudication is required before the reference can be treated
as validated ground truth. Despite a field name containing ``independent
benchmark'', the scores are rule-generated synthetic anchors; the label
does not make them empirically independent or human-validated.

\subsection{Analysis Plan}\label{analysis-plan}

Three analyses were conducted using existing workbook values only.
First, a structural audit checked row counts, unique identifiers, role
coverage, pattern counts, annotation completeness, and benchmark status.
Second, the stored composite was independently recomputed from the
documented formula and the workbook's EC, ER, CA, CP, and ET columns;
absolute differences were summarized. Third, an exploratory sensitivity
analysis recomputed scores under three named alternative weight
configurations and 10,000 random weight draws under each of two
Dirichlet schemes. The random generator seed was fixed at 20261027. Rank
stability was measured with Spearman's rank correlation relative to the
proposed weights.

The sensitivity analysis is a mathematical analysis of provisional
synthetic component values, not a test of predictive accuracy or
construct validity. The source files available for this revision did not
include a complete transformation script and per-case output table
sufficient to reproduce the earlier described three-condition
experiment. Therefore, the earlier condition means and associated
rating-error results are not reported in this version. No new LLM/API
calls, human annotations, or real employee records were used.

\section{Results}\label{results}

\subsection{Dataset and Formula
Audit}\label{dataset-and-formula-audit}

The structural audit identified 100 cases with 100 unique IDs, no
duplicate case IDs, and 10 roles represented by 10 cases each. All 100
benchmark records remained provisional. The annotation workbook
contained 500 criterion-level rows, with zero rows containing any of the
four core rating/coding fields. The stored ERA composite had a mean of
0.80874 and ranged from 0.667 to 0.947; these are descriptive summaries
of the supplied synthetic records, not model performance results.

\begin{longtable}[]{@{}
  >{\raggedright\arraybackslash}p{(\columnwidth - 4\tabcolsep) * \real{0.3333}}
  >{\raggedright\arraybackslash}p{(\columnwidth - 4\tabcolsep) * \real{0.3333}}
  >{\raggedright\arraybackslash}p{(\columnwidth - 4\tabcolsep) * \real{0.3333}}@{}}
\toprule\noalign{}
\begin{minipage}[b]{\linewidth}\raggedright
\textbf{Audit item}
\end{minipage} & \begin{minipage}[b]{\linewidth}\raggedright
\textbf{Result}
\end{minipage} & \begin{minipage}[b]{\linewidth}\raggedright
\textbf{Interpretation}
\end{minipage} \\
\midrule\noalign{}
\endhead
\bottomrule\noalign{}
\endlastfoot
Synthetic cases / unique IDs & 100 / 100 & No duplicate case IDs
found. \\
Roles & 10; 10 cases per role & Balanced by construction. \\
Criterion annotation rows & 500 & Five criterion rows per case. \\
Rows with any core annotation field completed & 0 of 500 & No human
annotations are present. \\
Provisional benchmark status & 100 of 100 cases & Reference scores are
not validated ground truth. \\
Mean absolute difference: stored vs recomputed ERA & 0.000265 &
Consistent with rounding in the workbook. \\
Maximum absolute difference & 0.000600 & No row exceeded the 0.001
tolerance. \\
Rows with difference \textgreater{} 0.001 & 0 & Arithmetic consistency
check passed at the stated tolerance. \\
\end{longtable}

The formula audit supports a narrow conclusion: the stored composite
values are arithmetically consistent with the documented equation to
within the specified tolerance. It does not validate the component
annotations, synthetic reference scores, weights, or substantive meaning
of the score.

\subsection{Exploratory Weight
Sensitivity}\label{exploratory-weight-sensitivity}

Under the proposed weights, the mean composite recalculated from the
available component fields was 0.808754. Three named alternative weight
sets produced Spearman rank correlations of 0.9446 (equal weights),
0.9685 (evidence-heavy), and 0.9552 (traceability-heavy) relative to the
proposed weights. In the 10,000-draw local Dirichlet analysis centered
on the proposed weights, the median rank correlation was 0.9641
(5th--95th percentile: 0.8872--0.9916), and 92.47\% of draws had $\rho \geq 0.90$. Under broad Dirichlet(1,1,1,1,1) perturbations, the median was
0.7563 (5th--95th percentile: 0.4797--0.9360), and 12.09\% of draws had $\rho \geq 0.90$.

\begin{longtable}[]{@{}
  >{\raggedright\arraybackslash}p{(\columnwidth - 4\tabcolsep) * \real{0.3333}}
  >{\raggedright\arraybackslash}p{(\columnwidth - 4\tabcolsep) * \real{0.3333}}
  >{\raggedright\arraybackslash}p{(\columnwidth - 4\tabcolsep) * \real{0.3333}}@{}}
\toprule\noalign{}
\begin{minipage}[b]{\linewidth}\raggedright
\textbf{Weight scheme}
\end{minipage} & \begin{minipage}[b]{\linewidth}\raggedright
\textbf{Mean composite}
\end{minipage} & \begin{minipage}[b]{\linewidth}\raggedright
\textbf{Spearman $\rho$ vs proposed}
\end{minipage} \\
\midrule\noalign{}
\endhead
\bottomrule\noalign{}
\endlastfoot
Proposed: 0.25, 0.25, 0.25, 0.15, 0.10 & 0.808754 & 1.0000 \\
Equal: 0.20 each & 0.805490 & 0.9446 \\
Evidence-heavy: 0.30, 0.30, 0.15, 0.15, 0.10 & 0.801874 & 0.9685 \\
Traceability-heavy: 0.20, 0.20, 0.20, 0.25, 0.15 & 0.805065 & 0.9552 \\
\end{longtable}

\begin{longtable}[]{@{}
  >{\raggedright\arraybackslash}p{(\columnwidth - 10\tabcolsep) * \real{0.1667}}
  >{\raggedright\arraybackslash}p{(\columnwidth - 10\tabcolsep) * \real{0.1667}}
  >{\raggedright\arraybackslash}p{(\columnwidth - 10\tabcolsep) * \real{0.1667}}
  >{\raggedright\arraybackslash}p{(\columnwidth - 10\tabcolsep) * \real{0.1667}}
  >{\raggedright\arraybackslash}p{(\columnwidth - 10\tabcolsep) * \real{0.1667}}
  >{\raggedright\arraybackslash}p{(\columnwidth - 10\tabcolsep) * \real{0.1667}}@{}}
\toprule\noalign{}
\begin{minipage}[b]{\linewidth}\raggedright
\textbf{Random-weight scheme}
\end{minipage} & \begin{minipage}[b]{\linewidth}\raggedright
\textbf{Draws}
\end{minipage} & \begin{minipage}[b]{\linewidth}\raggedright
\textbf{5th percentile $\rho$}
\end{minipage} & \begin{minipage}[b]{\linewidth}\raggedright
\textbf{Median $\rho$}
\end{minipage} & \begin{minipage}[b]{\linewidth}\raggedright
\textbf{95th percentile $\rho$}
\end{minipage} & \begin{minipage}[b]{\linewidth}\raggedright
\textbf{Share with $\rho \geq 0.90$}
\end{minipage} \\
\midrule\noalign{}
\endhead
\bottomrule\noalign{}
\endlastfoot
Local: Dirichlet(40 × proposed weights) & 10,000 & 0.8872 & 0.9641 &
0.9916 & 92.47\% \\
Broad: Dirichlet(1,1,1,1,1) & 10,000 & 0.4797 & 0.7563 & 0.9360 &
12.09\% \\
\end{longtable}

The sensitivity analysis indicates that rankings of these synthetic
cases are more stable when weights remain near the proposed
configuration and less stable when weights vary broadly. These
proportions describe analyst-selected mathematical perturbations; they
are not probabilities that the weights are correct. Because the
underlying component values are synthetic and provisional, the analysis
does not establish optimal weights, real-world rank stability, or
validity of employee ratings.

\section{Discussion}\label{discussion}

\subsection{What the Current Analysis
Establishes}\label{what-the-current-analysis-establishes}

The available evidence supports three limited conclusions. First, the
ERA equation is explicitly defined and can be recalculated from the
supplied normalized component fields. Second, the workbook has a
balanced synthetic structure across ten named roles, while its 500
annotation rows remain unfilled and its reference status remains
provisional. Third, case rankings based on these synthetic components
are relatively stable under local weight perturbations but more
sensitive under broad perturbations. These conclusions are reproducible
from the included workbook and script.

\subsection{What It Does Not
Establish}\label{what-it-does-not-establish}

The analysis does not show that ERA identifies correct
employee-performance ratings, that an AI system generated the listed
illustrative ratings, that a human reviewer agrees with the benchmark,
or that ERA improves decisions. The workbook's
``illustrative\_ai\_rating'' field is not accompanied by a model
identifier, prompt, run log, or raw model output; it is therefore
treated as an illustrative synthetic field rather than a reproducible AI
result. The synthetic reference scores are not independent human ground
truth. No human inter-rater reliability, explanation faithfulness,
fairness, reviewer usefulness, or organizational deployment result is
available.

\subsection{Implications for HRMS
Design}\label{implications-for-hrms-design}

ERA is best understood as a candidate audit layer that makes evidence
coverage, relevance, agreement, contradiction, and traceability
explicit. In an HRMS, component-level visibility may help a reviewer
identify why a composite score is high or low, but the score should
never conceal an undefined component or a critical contradiction. A
system should display the component profile and the underlying records,
allow abstention when evidence is insufficient, and preserve a contest
path. These are design implications, not claims that such benefits have
been demonstrated.

\subsection{Research Questions
Revisited}\label{research-questions-revisited}

RQ1 (operationalization) is addressed at the specification and
arithmetic-audit level. RQ2 (response to controlled evidence changes) is
not claimed as a finding in this version because the available
condition-level transformation artifacts were insufficient to reproduce
the earlier reported experiment. RQ3 (incremental value over simpler
baselines) remains unanswered because no independent outcome is
available. RQ4 (weight sensitivity) is addressed only as a mathematical
sensitivity analysis on synthetic components. RQ5 (reliability,
usefulness, fairness, and external validity) remains outside the
evidence currently available. These boundaries are part of the result,
not a request that the current paper be supplemented with invented data.

\section{Validity, Ethics, and
Limitations}\label{validity-ethics-and-limitations}

Construct validity is not established: the component definitions are a
proposal and have not been validated by a representative expert panel or
against independent outcomes. Criterion validity is not established
because the reference scores are synthetic and provisional. Internal
validity for an evidence-intervention effect cannot be assessed without
reproducible condition transformations and outputs. External validity is
not established because no real employee records or organization were
studied. Computational reproducibility is supported for the formula
audit and weight sensitivity analysis; human annotation reliability was
not measured.

The dataset is synthetic, with 100 constructed cases across 10 roles.
Its pattern frequencies are designed and do not estimate real-world
prevalence. All 500 core annotation records are blank. The workbook
includes rating-like strings that were capped at the stated 5-point
scale in a quality-corrected copy; the correction log preserves the
edits and raw evidence strings. This correction improves scale
consistency but does not validate the content. The illustrative rating
field lacks documented model provenance. Weight sensitivity depends on
analyst-selected perturbation distributions and provisional inputs.

No real employee data, personally identifiable information, human
participant data, or external LLM/API calls were used in the current
computational audit. This statement does not assert that institutional
ethics approval or exemption has been granted. Any later study involving
human participants or real employee data must follow the relevant
institutional and legal review, data minimization, access control,
retention, and governance requirements. ERA must not automatically
determine employment outcomes; meaningful review and contestability are
necessary safeguards.

\section{Conclusion}\label{conclusion}

This paper presents Evidence--Rating Alignment (ERA) as a proposed
measurement layer for auditing whether employee-performance ratings are
supported by criterion-relevant evidence. It defines five normalized
components---coverage, relevance, criterion agreement, traceability, and
contradiction penalty---and specifies how they can be embedded in an
auditable HRMS workflow. An audit of the available 100-case synthetic
workbook found 100 unique cases across 10 roles, 500 uncompleted
criterion-level annotation rows, provisional benchmark status for every
case, and stored composite values that matched the documented equation
within a maximum absolute discrepancy of 0.0006. An exploratory
sensitivity analysis showed stronger rank stability under local than
broad weight perturbations on these synthetic values.

The contribution is a framework specification and reproducible
computational audit, not a validated employee-performance instrument. No
human or LLM results are claimed, and no real-world accuracy, fairness,
or decision-improvement effect is inferred. The paper is complete within
that declared scope: it reports what the existing artifacts can support
and explicitly excludes claims that they cannot support. ERA should
remain a research prototype and must not be used as a standalone
employment decision score.

\section{Reproducibility Statement}\label{reproducibility-statement}

The analysis uses the existing QC-corrected workbook, the ERA formula
stated in Section 5, and a fixed-seed script that performs structural
checks, recomputes the composite, evaluates four named weight
configurations, and samples 10,000 weight vectors for each of two
sensitivity schemes. Machine-readable audit outputs and the script are
included in the accompanying reproducibility package. All computed
results are descriptive of synthetic data and should be interpreted
within the limitations stated above.

\begin{thebibliography}{23}
\bibitem{r1} A. Bujold et al., ``Responsible artificial intelligence in human
resources management: A review of the empirical literature,'' AI and
Ethics, vol. 4, pp. 1185--1200, 2024, doi: 10.1007/s43681-023-00325-1.
\bibitem{r2} A. Varma, C. Dawkins, and K. Chaudhuri, ``Artificial
intelligence and people management: A critical assessment through the
ethical lens,'' Human Resource Management Review, vol. 33, no. 1,
100923, 2023, doi: 10.1016/j.hrmr.2022.100923.
\bibitem{r3} T. Miller, ``Explanation in artificial intelligence: Insights
from the social sciences,'' Artificial Intelligence, vol. 267, pp.
1--38, 2019, doi: 10.1016/j.artint.2018.07.007.
\bibitem{r4} M. T. Ribeiro, S. Singh, and C. Guestrin, ``Why should I trust
you? Explaining the predictions of any classifier,'' in Proc. KDD, 2016,
pp. 1135--1144, doi: 10.1145/2939672.2939778.
\bibitem{r5} A. Köchling and M. C. Wehner, ``Discriminated by an algorithm: A
systematic review of discrimination and fairness by algorithmic
decision-making in HR recruitment and development,'' Business Research,
vol. 13, pp. 795--848, 2020, doi: 10.1007/s40685-020-00134-w.
\bibitem{r6} D. T. Newman, N. J. Fast, and D. J. Harmon, ``When eliminating
bias isn't fair: Algorithmic reductionism and procedural justice in
human resource decisions,'' Organizational Behavior and Human Decision
Processes, vol. 160, pp. 149--167, 2020, doi:
10.1016/j.obhdp.2020.03.008.
\bibitem{r7} J. Qin, Y. Jia, X. Luo, Q. V. Liao, and Z. Huang, ``Perceived
fairness of human managers compared with artificial intelligence in
employee performance evaluation,'' Journal of Management Information
Systems, vol. 40, no. 4, pp. 1039--1070, 2023, doi:
10.1080/07421222.2023.2267316.
\bibitem{r8} K. C. Kellogg, M. A. Valentine, and A. Christin, ``Algorithms at
work: The new contested terrain of control,'' Academy of Management
Annals, vol. 14, no. 1, pp. 366--410, 2020, doi:
10.5465/annals.2018.0174.
\bibitem{r9} A. Keegan and J. Meijerink, ``Algorithmic management in
organizations? From edge case to center stage,'' Annual Review of
Organizational Psychology and Organizational Behavior, vol. 12, pp.
395--422, 2025, doi: 10.1146/annurev-orgpsych-110622-070928.
\bibitem{r10} N. Li, H. Zhou, and M. Xu, ``From text to insight: Leveraging
large language models for performance evaluation in management,''
Personnel Psychology, vol. 79, no. 2, pp. 251--267, 2026, doi:
10.1111/peps.70020.
\bibitem{r11} P. Tambe, P. Cappelli, and V. Yakubovich, ``Artificial
intelligence in human resources management: Challenges and a path
forward,'' California Management Review, vol. 61, no. 4, pp. 15--42,
2019, doi: 10.1177/0008125619867910.
\bibitem{r12} A. Malik, P. Budhwar, and B. A. Kazmi, ``Artificial
intelligence (AI)-assisted HRM: Towards an extended strategic
framework,'' Human Resource Management Review, vol. 33, no. 1, 100940,
2023, doi: 10.1016/j.hrmr.2022.100940.
\bibitem{r13} N. Noponen et al., ``Taylorism on steroids or enabling
autonomy? A systematic review of algorithmic management,'' Management
Review Quarterly, vol. 74, pp. 1695--1721, 2024, doi:
10.1007/s11301-023-00345-5.
\bibitem{r14} S. Bankins, P. Formosa, Y. Griep, and D. Richards, ``AI
decision making with dignity? Contrasting workers' justice perceptions
of human and AI decision making in a human resource management
context,'' Information Systems Frontiers, vol. 24, pp. 857--875, 2022,
doi: 10.1007/s10796-021-10223-8.
\bibitem{r15} A. Shulner-Tal et al., ``Who made that decision and why? Users'
perceptions of human versus AI decision-making and the power of
explainable AI,'' International Journal of Human--Computer Interaction,
2024, doi: 10.1080/10447318.2024.2348843.
\bibitem{r16} X. Xu, Y. Chen, and J. Dong, ``Shaping the fairness journey:
The roles of AI literacy, explanation, and interpersonal interaction in
AI interviews,'' International Journal of Human-Computer Studies, vol.
205, 103629, 2025, doi: 10.1016/j.ijhcs.2025.103629.
\bibitem{r17} S. Amershi et al., ``Guidelines for human-AI interaction,'' in
Proc. CHI, 2019, doi: 10.1145/3290605.3300233.
\bibitem{r18} N. Bennett and C. L. Martin, ``AI as a talent management tool:
An organizational justice perspective,'' Business Horizons, vol. 69, no.
2, pp. 241--252, 2026, doi: 10.1016/j.bushor.2025.03.005.
\bibitem{r19} M. J. Page et al., ``The PRISMA 2020 statement: an updated
guideline for reporting systematic reviews,'' BMJ, vol. 372, n71, 2021,
doi: 10.1136/bmj.n71.
\bibitem{r20} I. Kekez, L. Lauwaert, and N. Begičević Ređep, ``Is artificial
intelligence (AI) research biased and conceptually vague? A systematic
review of research on bias and discrimination in the context of using AI
in human resource management,'' Technology in Society, vol. 81, 102818,
2025, doi: 10.1016/j.techsoc.2025.102818.
\bibitem{r21} OECD/European Union/European Commission Joint Research Centre,
Handbook on Constructing Composite Indicators: Methodology and User
Guide. Paris: OECD Publishing, 2008, doi: 10.1787/9789264043466-en.
\bibitem{r22} W. Becker, M. Saisana, P. Paruolo, and I. Vandecasteele,
``Weights and importance in composite indicators: Closing the gap,''
Ecological Indicators, vol. 80, pp. 12--22, 2017, doi:
10.1016/j.ecolind.2017.03.056.
\bibitem{r23} R. F. Naoum, T. Szakadáti, and G. Balogh, ``Artificial
Intelligence (AI) in human resource management (HRM): a systematic
review of its dual impact on diversity, equity, and inclusion (DEI),''
Management Review Quarterly, published online 8 January 2026, doi:
10.1007/s11301-025-00580-y.
\end{thebibliography}
\end{document}
